\begin{spacing}{1.5}PixelLab provides a practical environment for studying digital image processing through a local desktop interface and an offline Android application. The project addresses the effort required to move between course references, independent scripts, parameter settings and result interpretation. A common registry describes operations and validates their inputs, while the interface presents processed images, operation-specific figures, measurements and Python examples. The implemented system contains 59 registered operations and a catalog of 86 course topics covering fundamentals, enhancement, color processing, transforms, segmentation, restoration, compression and morphology. Android executes bundled Python algorithms through Chaquopy without requiring a running computer or processing server. Evaluation combines fresh bounded execution checks, exact classroom examples, lossless analysis-image round trips and retained platform verification records. All 59 default operation paths produced nonempty finite uint8 results on the report sample. The normalized equalization example maps intensities 52 and 55 to 0 and 12, and all four teaching lossless codecs recover their grayscale analysis input exactly. Historical records separately document 47 desktop tests and 272 Android integration checks. These findings support the implemented workflow and numerical demonstrations. Evaluation on physical ARM64 phones and measurement of student learning outcomes remain future work.\par\medskip\textbf{Keywords:} digital image processing, offline laboratory, Android, course alignment, numerical verification.\end{spacing}